\section{相关工作}

\paragraph{由 LLM/VLM 驱动的 GUI 智能体。}
LLM 和 VLM 的出现推动了 GUI 智能体的发展，使其能够理解自然语言指令，并在移动端~\citep{rawles2024androidworld}、Web~\citep{zhou2023webarena, deng2023mind2web} 和桌面环境~\citep{xie2024osworld, zhang2024ufo, zhang2025ufo2} 中执行复杂任务。
早期研究关注自主智能体框架~\citep{zheng2024gpt, wu2024copilot, jia2024agentstore}，这些框架提示商业模型通过代码生成~\citep{sun2024survey, sun2025scienceboard} 或工具调用~\citep{zhang2025appagent,zhang2025api} 与操作系统交互。
随着开源且可定制智能体的需求上升，另一条工作线聚焦于训练 LLM/VLM 来增强其智能体能力，包括 GUI 理解、规划和执行~\citep{xu2024aguvis, qin2025ui, chen2024guicourse, chai2024amex}。这些工作的关键在于收集 GUI 专用训练数据，例如 OCR 标注~\citep{lu2024gui}、界面摘要~\citep{you2024ferret}、问答对~\citep{chen2024guicourse} 和大规模任务演示~\citep{deng2023mind2web, rawles2023androidinthewild, zhang2024android, zheng2025vem, sun2024genesis, cheng2024vision, zhang2025breaking}。

智能体开发的一个核心需求，是能够与部署在虚拟机和基于 Chrome 的浏览器中的真实 GUI 环境交互。早期智能体操作 HTML 或可访问性树等结构化元数据~\citep{zhou2023webarena, he2024webvoyager}，但这类表示在不同平台上脆弱且不一致~\citep{cheng2024seeclick, xu2024aguvis}。因此，近期趋势转向\emph{以视觉为中心的范式}，即智能体像人一样，通过原始截图感知界面，并使用鼠标和键盘输入进行交互~\citep{openai2025operator, anthropic2024cuda}。在这一设定下，一个核心挑战随之出现：把自然语言指令定位到具体 GUI 区域，本文称之为\textit{GUI 视觉定位}。

\paragraph{GUI 视觉定位。}
给定一张 GUI 截图和一条自然语言指令，GUI 视觉定位旨在找到用于交互的目标区域。
虽然这一任务在概念上与自然图像中的定位相关，但 GUI 布局具有高语义密度和结构规律性，因此带来独特挑战~\citep{cheng2024seeclick, gou2024navigating}。一种常见做法是把 GUI 定位表述为基于文本的坐标预测任务，让模型把点位置（\eg, \texttt{x=..., y=...}）作为语言 token 输出~\citep{chen2023shikra, wang2024qwen2}。由于简单且兼容现有 LLM/VLM，这种表述被广泛采用。为了提升性能，已有工作同时扩展模型和训练数据~\citep{lin2024showui, yang2024aria, pawlowski2024tinyclick, tang2025think, bai2025qwen2, qin2025ui, xie2025scaling}。UGround~\citep{gou2024navigating} 提出一种用于合成多样 GUI 定位样本的数据流水线，而 OS-Atlas~\citep{wu2024atlas} 提供了多平台数据集和统一的 GUI 动作模型。最近，\citet{xu2025attention} 提出了一种无需训练的方法，通过利用 VLM 的内部注意力来执行 GUI 定位。

尽管这些方法取得了成功，基于坐标的方法仍存在关键限制，包括空间归纳偏置较弱、点监督存在歧义，以及视觉特征与动作目标之间存在分辨率不匹配。
本文提出了对主流坐标方法的一种有力替代方案：GUI-Actor，一种面向 GUI 智能体的新型\textit{无坐标定位框架}。
它引入一个 \texttt{<ACTOR>} token，并通过注意力式动作头直接关注相关图像 patch，从而实现更接近人类的定位方式，同时缓解基于坐标方法的局限。
